\documentclass{ximera}

\input{../../../preamble.tex}


\definecolor{fvdnavy}{HTML}{071D43}
\definecolor{fvdcyan}{HTML}{20BFC6}
\definecolor{fvdcyanlight}{HTML}{EFFCFB}
\definecolor{fvdmagenta}{HTML}{F10D69}
\definecolor{fvdmagentalight}{HTML}{FFF1F7}
\definecolor{fvdtext}{HTML}{163044}





%=================================================
% Pedagogische omgevingen
% PDF: tcolorbox
% HTML: learning-box
%=================================================

\newenvironment{voorkennis}
{%
    \ifdefined\HCode
        \HCode{%
<section class="learning-box learning-box--prior">
<div class="learning-box__header">
<span class="learning-box__icon" aria-hidden="true"></span>
<span class="learning-box__title">Voorkennis</span>
</div>
<div class="learning-box__content">
        }%
        \begin{itemize}
    \else
       \begin{tcolorbox}[
    enhanced,
    breakable,
    colback=fvdcyanlight,
    colframe=fvdcyan,
    coltext=fvdtext,
    boxrule=0.5pt,
    leftrule=4pt,
    arc=4mm,
    outer arc=4mm,
    left=7mm,
    right=6mm,
    top=5mm,
    bottom=5mm,
    before skip=6mm,
    after skip=6mm,
    title=Voorkennis,
    coltitle=fvdcyan,
    fonttitle=\bfseries\Large,
    attach title to upper={\par\smallskip},
    boxed title style={
        empty,
        left=0mm,
        right=0mm,
        top=0mm,
        bottom=0mm
    },
    overlay={
        \draw[
            fvdcyan,
            line width=0.5pt
        ]
        ([xshift=42mm,yshift=-13mm]frame.north west)
        --
        ([xshift=-7mm,yshift=-13mm]frame.north east);
    }
]
        \begin{itemize}
    \fi
}
{%
    \end{itemize}
    \ifdefined\HCode
        \HCode{%
</div>
</section>
        }%
    \else
        \end{tcolorbox}
    \fi
}


\newenvironment{leerdoelen}
{%
    \ifdefined\HCode
        \HCode{%
<section class="learning-box learning-box--goals">
<div class="learning-box__header">
<span class="learning-box__icon" aria-hidden="true"></span>
<span class="learning-box__title">Leerdoelen</span>
</div>
<div class="learning-box__content">
        }%
        \begin{itemize}
    \else
\begin{tcolorbox}[
    enhanced,
    breakable,
    colback=fvdmagentalight,
    colframe=fvdmagenta,
    coltext=fvdtext,
    boxrule=0.5pt,
    leftrule=4pt,
    arc=4mm,
    outer arc=4mm,
    left=7mm,
    right=6mm,
    top=5mm,
    bottom=5mm,
    before skip=6mm,
    after skip=6mm,
    title=Leerdoelen,
    coltitle=fvdmagenta,
    fonttitle=\bfseries\Large,
    attach title to upper={\par\smallskip},
    boxed title style={
        empty,
        left=0mm,
        right=0mm,
        top=0mm,
        bottom=0mm
    },
    overlay={
        \draw[
            fvdmagenta,
            line width=0.5pt
        ]
        ([xshift=42mm,yshift=-13mm]frame.north west)
        --
        ([xshift=-7mm,yshift=-13mm]frame.north east);
    }
]
        \begin{itemize}
    \fi
}
{%
    \end{itemize}
    \ifdefined\HCode
        \HCode{%
</div>
</section>
        }%
    \else
        \end{tcolorbox}
    \fi
}


\addPrintStyle{../../..}

\title{Bewegen onder invloed van velden}
\author{Ingmar Herreman}

\begin{abstract}
In de vorige hoofdstukken leerden we bewegingen beschrijven en onderzochten
we hoe een resulterende kracht volgens de wetten van Newton een beweging
verandert. In dit hoofdstuk brengen we die ideeën samen met een begrip dat
je al eerder ontmoette: het veld.

We starten bij gravitatie. Newtons gravitatiewet en het gravitatieveld
worden kort geactiveerd en vervolgens gekoppeld aan de cirkelbeweging.
Zo ontdekken we waarom een satelliet voortdurend naar de aarde kan vallen
zonder het aardoppervlak te bereiken. Uit de combinatie van gravitatie,
Newton II en centripetale versnelling leiden we de baansnelheid en de
omlooptijd van een satelliet af. Daarbij verschijnt ook het verband van
Kepler op een natuurlijke manier.

Daarna maken we transfer naar elektrische en magnetische velden. Een
geladen deeltje in een homogeen elektrisch veld ondervindt een constante
kracht en vertoont daardoor een beweging die sterk lijkt op een worp in
een homogeen gravitatieveld. In een homogeen magnetisch veld kan de
magnetische kracht daarentegen voortdurend loodrecht op de snelheid staan
en zo een cirkelbaan veroorzaken.

Doorheen het hoofdstuk staat één gedachte centraal:
\[
\boxed{\text{veld}\longrightarrow\text{kracht}\longrightarrow
\text{versnelling}\longrightarrow\text{beweging}.}
\]
Dezelfde wetten van de mechanica verbinden zo vallende voorwerpen,
satellieten en bewegende geladen deeltjes.
\end{abstract}

\begin{document}

% FVD_CHAPTER_DATA_START
% Automatisch gegenereerd uit index.tex.
% Niet handmatig aanpassen.
\fvdchapterdata{1}{Beweging en kracht}{7}
% FVD_CHAPTER_DATA_END

\maketitle

% ============================================================
% 1 VOORKENNIS EN LEERDOELEN
% ============================================================

\begin{voorkennis}
  \item Je kan de wetten van Newton gebruiken, in het bijzonder
        $\sum\vec F=m\vec a$.
  \item Je kan een vrijlichaamsdiagram maken en krachten vectorieel analyseren.
  \item Je kent bij een eenparige cirkelbeweging
        $v=\omega r$ en $a_c=v^2/r=\omega^2r$.
  \item Je kent uit eerdere leerstof de begrippen gravitatieveld,
        elektrisch veld en magnetisch veld.
  \item Je kent de gravitatiekracht, elektrische kracht en magnetische
        kracht op een bewegende lading.
  \item Je kan bewegingen in één en twee dimensies analyseren.
  \item Je kan evenredigheden en machtsverbanden interpreteren.
\end{voorkennis}

\begin{leerdoelen}
  \item Je kan Newtons gravitatiewet en het gravitatieveld analyseren en
        kwantificeren.
  \item Je kan het onderscheid maken tussen afstand tot een
        planeetmiddelpunt, baanstraal en hoogte boven het oppervlak.
  \item Je kan het verband tussen gravitatiekracht en centripetale
        versnelling van een satelliet verklaren en gebruiken.
  \item Je kan voor een cirkelvormige satellietbaan de baansnelheid en
        omlooptijd afleiden, interpreteren en berekenen.
  \item Je kan het verband $T^2\sim r^3$ analyseren.
  \item Je kan uitleggen waarom astronauten in een baan om de aarde
        gewichtloos lijken zonder dat de gravitatie daar nul is.
  \item Je kan de beweging van een geladen deeltje in een homogeen
        elektrisch veld analyseren.
  \item Je kan de magnetische kracht op een bewegende lading koppelen aan
        een cirkelbeweging wanneer $\vec v\perp\vec B$.
  \item Je kan overeenkomsten en verschillen tussen gravitatie-, elektrische
        en magnetische krachtwerking analyseren.
  \item Je kan nieuwe veldproblemen systematisch mathematiseren vanuit
        veld, kracht, Newton II en kinematica.
\end{leerdoelen}


% ============================================================
% 2 OPENING
% ============================================================

\FVDsection{Eén mechanica, verschillende velden}

Een appel valt naar de aarde. Een kunstmaan beweegt rond de aarde.
Een elektron wordt versneld tussen twee geladen platen. Een ion buigt
af in een magnetisch veld.

Op het eerste gezicht lijken dit vier totaal verschillende verschijnselen.
Toch kunnen we ze met dezelfde denkwijze onderzoeken.

Een veld bepaalt welke kracht een geschikt systeem in een bepaald punt
ondervindt. Die kracht bepaalt volgens Newton II de versnelling, en de
versnelling bepaalt hoe de snelheid verandert:

\[
\boxed{
\text{veld}
\longrightarrow
\text{kracht}
\longrightarrow
\text{versnelling}
\longrightarrow
\text{beweging}.
}
\]

Dit hoofdstuk is daarom geen verzameling nieuwe formules. Het is een
synthese van begrippen die we eerder afzonderlijk leerden gebruiken.


\begin{oefening}[Vier bewegingen, één denkwijze]{\oefU\ESS}

Beschouw de volgende situaties:

\begin{enumerate}
  \item een steen valt in vacuüm nabij het aardoppervlak;
  \item een satelliet beweegt in een cirkelbaan rond de aarde;
  \item een positief geladen deeltje vertrekt uit rust in een homogeen
        elektrisch veld;
  \item een positief geladen deeltje beweegt loodrecht op een homogeen
        magnetisch veld.
\end{enumerate}

\begin{question}
Geef voor elke situatie aan welke kracht de beweging beïnvloedt.
\end{question}

\begin{question}
Voorspel voor elke situatie of de grootte van de snelheid constant kan
blijven.
\end{question}

\begin{question}
Voorspel voor elke situatie of de richting van de snelheid verandert.
\end{question}

\begin{question}
Welke twee situaties verwacht je met een cirkelbeweging te kunnen
verbinden?
\end{question}

Bewaar je antwoorden. Aan het einde van het hoofdstuk keren we naar deze
vergelijking terug.

\end{oefening}


% ============================================================
% DEEL A GRAVITATIE
% ============================================================

\FVDsection{Gravitatie opnieuw geactiveerd}

Alle massa's trekken elkaar gravitatieel aan. Voor twee puntmassa's
$m_1$ en $m_2$ op onderlinge afstand $r$ geldt voor de grootte van de
gravitatiekracht:

\[
\boxed{
F_g=G\frac{m_1m_2}{r^2}.
}
\]

Hierbij is

\[
G\approx 6{,}67\cdot10^{-11}\,
\mathrm{\frac{N\,m^2}{kg^2}}
\]

de gravitatieconstante.

Voor een bolsymmetrisch hemellichaam mogen we, buiten het hemellichaam,
de massa voor de gravitatieberekening behandelen alsof ze in het
middelpunt geconcentreerd is.

\begin{opmerking}
De afstand $r$ in de gravitatiewet is dus niet de hoogte boven het
oppervlak. Ze is de afstand tussen de massamiddelpunten.
\end{opmerking}


\FVDsubsection{Een inverse-kwadratenwet}

Als de massa's constant blijven:

\[
F_g\sim\frac1{r^2}.
\]

Dus:

\[
r\rightarrow 2r
\quad\Longrightarrow\quad
F_g\rightarrow\frac14F_g,
\]

en:

\[
r\rightarrow 3r
\quad\Longrightarrow\quad
F_g\rightarrow\frac19F_g.
\]

Het verband is niet lineair.


\begin{oefening}[Redeneren met de gravitatiewet]{\oefB\ESS}

Twee puntmassa's oefenen een gravitatiekracht $F$ op elkaar uit.

\begin{question}
Hoe groot wordt de kracht als één massa verdubbelt en de andere
grootheden gelijk blijven?
\end{question}

\begin{question}
Hoe groot wordt de kracht als beide massa's verdubbelen?
\end{question}

\begin{question}
Hoe groot wordt de kracht als de afstand tussen de massa's verdubbelt?
\end{question}

\begin{question}
Met welke factor moet de afstand veranderen om de kracht viermaal zo
klein te maken?
\end{question}

\end{oefening}


% ============================================================
% 4 GRAVITATIEVELD
% ============================================================

\FVDsection{Het gravitatieveld}

In plaats van voor elke mogelijke proefmassa opnieuw de gravitatiekracht
te berekenen, beschrijven we de ruimte rond een massa met een veld.

\begin{definitie}
De \textbf{gravitatieveldsterkte} $\vec g$ in een punt is de
gravitatiekracht per eenheid proefmassa:

\[
\boxed{
\vec g=\frac{\vec F_g}{m}.
}
\]
\end{definitie}

Voor een bolsymmetrische bronmassa $M$:

\[
F_g=G\frac{Mm}{r^2}.
\]

Delen door $m$ geeft:

\[
\boxed{
g(r)=G\frac{M}{r^2}.
}
\]

De vector $\vec g$ wijst naar het middelpunt van de aantrekkende massa.

De eenheden

\[
\mathrm{\frac{N}{kg}}
\]

en

\[
\mathrm{\frac{m}{s^2}}
\]

zijn voor de gravitatieveldsterkte equivalent.


\FVDsubsection{Waarom is $g$ nabij de aarde ongeveer constant?}

Op het aardoppervlak:

\[
g_0=G\frac{M_A}{R_A^2}.
\]

Op hoogte $h$ boven het aardoppervlak:

\[
r=R_A+h
\]

en dus:

\[
\boxed{
g(h)=G\frac{M_A}{(R_A+h)^2}.
}
\]

Voor hoogtes die zeer klein zijn tegenover $R_A$ verandert $r$ relatief
weinig. Daarom mogen we in veel val- en worpproblemen nabij het
aardoppervlak $g$ als ongeveer constant beschouwen.

Voor satellieten op grote hoogte mag dat niet zomaar.


\begin{oefening}[Hoogte is geen baanstraal]{\oefB\ESS}

Een satelliet bevindt zich op hoogte $h$ boven een planeet met straal
$R$.

\begin{question}
Welke afstand moet in
\[
g=G\frac{M}{r^2}
\]
worden gebruikt?
\end{question}

\begin{question}
Schrijf $g$ als functie van $R$ en $h$.
\end{question}

\begin{question}
Een satelliet bevindt zich op hoogte $h=R$. Hoe groot is daar de
gravitatieveldsterkte als fractie van de veldsterkte aan het oppervlak?
\end{question}

\end{oefening}


\begin{oefening}[Een grafiek van $g(r)$]{\oefU\ESS}

Voor punten buiten een bolsymmetrische planeet geldt:

\[
g(r)=G\frac{M}{r^2}.
\]

\begin{question}
Beschrijf het functietype en het gedrag van $g(r)$ voor $r>R$.
\end{question}

\begin{question}
Wat gebeurt er met $g$ als $r$ zeer groot wordt?
\end{question}

\begin{question}
Is $g(r)$ lineair dalend? Motiveer wiskundig.
\end{question}

\begin{question}
Welke grafiek zou je kunnen maken om het inverse-kwadratenverband te
lineariseren?
\end{question}

\end{oefening}


% ============================================================
% 5 NEWTONS KANON
% ============================================================

\FVDsection{Van vallende steen naar satelliet}

Stel je een voorwerp voor dat vanaf een zeer hoge berg horizontaal wordt
weggeschoten.

Bij een kleine beginsnelheid valt het voorwerp naar de aarde en raakt
het oppervlak.

Bij een grotere beginsnelheid legt het tijdens zijn val een grotere
horizontale afstand af.

In het ideale gedachte-experiment kan de beginsnelheid zo groot worden
dat het aardoppervlak door zijn kromming even snel ``wegvalt'' als het
voorwerp naar de aarde valt. Het voorwerp blijft dan voortdurend vallen,
maar bereikt de aarde niet.

Dat is de essentie van een baanbeweging.

\begin{center}
\[
\boxed{
\text{een satelliet in een baan is voortdurend in vrije val}.
}
\]
\end{center}

De gravitatiekracht verdwijnt dus niet. Integendeel: ze is precies de
kracht die de baan voortdurend doet afbuigen.


\begin{oefening}[Een hardnekkige misconceptie]{\oefU\ESS}

Een leerling zegt:

\begin{quote}
``Een satelliet blijft in de ruimte omdat daar bijna geen zwaartekracht
meer is.''
\end{quote}

\begin{question}
Waarom kan die uitspraak niet verklaren waarom de satelliet een
cirkelbaan volgt?
\end{question}

\begin{question}
Wat zou er volgens Newton I gebeuren als de resulterende kracht werkelijk
nul was?
\end{question}

\begin{question}
Welke richting heeft de versnelling van een satelliet in een
cirkelbaan?
\end{question}

\begin{question}
Welke kracht kan die versnelling veroorzaken?
\end{question}

\end{oefening}


% ============================================================
% 6 SATELLIET CIRKELBAAN
% ============================================================

\FVDsection{Gravitatie als radiale resultante}

Beschouw een satelliet met massa $m$ in een cirkelbaan met straal $r$
rond een planeet met massa $M$.

In het ideale model verwaarlozen we andere krachten.

De enige kracht op de satelliet is dan de gravitatiekracht:

\[
F_g=G\frac{Mm}{r^2}.
\]

De satelliet voert een cirkelbeweging uit, dus:

\[
a_c=\frac{v^2}{r}.
\]

Volgens Newton II in de radiale richting:

\[
\sum F_r=ma_c.
\]

Omdat de gravitatiekracht de volledige radiale resultante levert:

\[
G\frac{Mm}{r^2}
=
m\frac{v^2}{r}.
\]

Dit is de centrale vergelijking voor een ideale cirkelvormige
satellietbaan.

\begin{opmerking}
We tekenen in het vrijlichaamsdiagram alleen $\vec F_g$. We tekenen
niet nog een extra ``centripetale kracht''. De gravitatiekracht
\emph{is} hier de radiale resultante.
\end{opmerking}


% ============================================================
% 7 BAANSNELHEID
% ============================================================

\FVDsection{De baansnelheid afleiden}

Vertrek van:

\[
G\frac{Mm}{r^2}
=
m\frac{v^2}{r}.
\]

De satellietmassa $m$ valt weg:

\[
G\frac{M}{r^2}
=
\frac{v^2}{r}.
\]

Vermenigvuldig met $r$:

\[
\frac{GM}{r}=v^2.
\]

Dus:

\[
\boxed{
v_{\mathrm{baan}}
=
\sqrt{\frac{GM}{r}}.
}
\]

Deze formule vertelt meer dan alleen hoe we een getal berekenen.

Bij een gegeven centrale massa $M$:

\[
\boxed{
v_{\mathrm{baan}}\sim\frac1{\sqrt r}.
}
\]

Een satelliet in een grotere cirkelbaan beweegt dus trager.


\FVDsubsection{Waarom verdwijnt de satellietmassa?}

Een zwaardere satelliet ondervindt een grotere gravitatiekracht:

\[
F_g\sim m.
\]

Maar volgens Newton II is ook meer kracht nodig om dezelfde versnelling
aan een grotere massa te geven:

\[
F=ma.
\]

De factor $m$ verschijnt aan beide kanten en valt weg.

Dat is dezelfde diepe eigenschap die we bij vrije val nabij het
aardoppervlak zagen.


\begin{oefening}[Baansnelheid analyseren]{\oefU\ESS}

Twee satellieten bewegen in cirkelbanen rond dezelfde planeet. Voor de
tweede satelliet geldt:

\[
r_2=4r_1.
\]

\begin{question}
Bepaal zonder getallen de verhouding $v_2/v_1$.
\end{question}

\begin{question}
Welke satelliet beweegt sneller?
\end{question}

\begin{question}
Hangt je antwoord af van de massa's van de satellieten?
\end{question}

\begin{question}
Leg uit waarom de uitspraak ``een hogere satelliet moet sneller bewegen
om niet te vallen'' fout is voor ideale cirkelbanen.
\end{question}

\end{oefening}


\begin{oefening}[Satelliet rond de aarde]{\oefB\ESS}

Gebruik:

\[
M_A=5{,}97\cdot10^{24}\,\mathrm{kg},
\qquad
R_A=6{,}37\cdot10^6\,\mathrm m.
\]

Een satelliet beweegt op hoogte

\[
h=400\,\mathrm{km}
\]

in een ideale cirkelbaan.

\begin{question}
Bepaal de baanstraal $r$.
\end{question}

\begin{question}
Bereken de baansnelheid.
\end{question}

\begin{question}
Vergelijk $g$ op die hoogte met $g$ aan het aardoppervlak.
\end{question}

\begin{question}
Leg op basis van je resultaat uit waarom ``geen zwaartekracht'' geen
goede verklaring is voor gewichtloosheid in een lage baan om de aarde.
\end{question}

\end{oefening}


% ============================================================
% 8 OMPLOOPTIJD
% ============================================================

\FVDsection{De omlooptijd}

Voor een eenparige cirkelbeweging is:

\[
v=\frac{2\pi r}{T}.
\]

Voor een satelliet in een ideale cirkelbaan vonden we:

\[
v=\sqrt{\frac{GM}{r}}.
\]

We stellen beide uitdrukkingen gelijk:

\[
\frac{2\pi r}{T}
=
\sqrt{\frac{GM}{r}}.
\]

Na herschikken:

\[
\boxed{
T
=
2\pi\sqrt{\frac{r^3}{GM}}.
}
\]

Voor eenzelfde centrale massa $M$:

\[
\boxed{
T\sim r^{3/2}.
}
\]

Een satelliet in een grotere baan heeft dus een langere omlooptijd.


\begin{oefening}[Verhoudingen zonder rekenmachine]{\oefU\ESS}

Twee satellieten bewegen rond dezelfde planeet in ideale cirkelbanen.

\[
r_2=4r_1.
\]

\begin{question}
Bepaal $T_2/T_1$.
\end{question}

\begin{question}
Bepaal opnieuw $v_2/v_1$.
\end{question}

\begin{question}
Leg uit waarom een grotere baan tegelijk een kleinere baansnelheid en
een grotere omlooptijd oplevert.
\end{question}

\end{oefening}


% ============================================================
% 9 KEPLER III
% ============================================================

\FVDsection{Keplers derde wet verschijnt uit Newton}

Kwadrateren van

\[
T
=
2\pi\sqrt{\frac{r^3}{GM}}
\]

geeft:

\[
\boxed{
T^2
=
\frac{4\pi^2}{GM}r^3.
}
\]

Voor objecten die rond dezelfde centrale massa bewegen, is

\[
\frac{4\pi^2}{GM}
\]

constant. Dus:

\[
\boxed{
T^2\sim r^3.
}
\]

Dit is de vorm van Keplers derde wet voor cirkelbanen.

Het verband is didactisch belangrijk omdat het laat zien dat een
empirisch patroon in planeetbewegingen verklaard kan worden vanuit
Newtons dynamica en gravitatiewet.


\FVDsubsection{Lineariseren}

Een grafiek van $T$ tegen $r$ is niet lineair.

Maar:

\[
T^2
=
\left(\frac{4\pi^2}{GM}\right)r^3.
\]

Een grafiek van $T^2$ tegen $r^3$ hoort dus een rechte door de oorsprong
te geven met richtingscoëfficiënt:

\[
\boxed{
a=\frac{4\pi^2}{GM}.
}
\]

Uit de helling kunnen we zelfs de massa van het centrale hemellichaam
bepalen:

\[
\boxed{
M=\frac{4\pi^2}{Ga}.
}
\]


\begin{oefening}[Een planeet wegen zonder weegschaal]{\oefG}

Van verschillende manen rond dezelfde planeet zijn de baanstraal $r$
en omlooptijd $T$ bekend.

\begin{question}
Welke twee grootheden zet je tegen elkaar uit om volgens het
cirkelbaanmodel een rechte te verwachten?
\end{question}

\begin{question}
Wat is de fysische betekenis van de helling van die rechte?
\end{question}

\begin{question}
Leid een formule af waarmee je uit de helling de massa van de planeet
kan bepalen.
\end{question}

\begin{question}
Waarom is dit een voorbeeld van indirect meten?
\end{question}

\begin{question}
Welke aannames van het model kunnen bij echte gegevens afwijkingen
veroorzaken?
\end{question}

\end{oefening}


% ============================================================
% 10 GEOSTATIONAIR
% ============================================================

\FVDsection{Een satelliet die boven hetzelfde punt blijft}

Communicatie- en weersatellieten kunnen zo worden geplaatst dat ze voor
een waarnemer op aarde boven ongeveer hetzelfde punt aan de hemel
blijven.

Zo'n satelliet noemen we \textbf{geostationair}.

Daarvoor volstaat niet alleen een bepaalde omlooptijd. De baan moet ook
geschikt georiënteerd zijn.

Voor een geostationaire satelliet geldt in het ideale model onder meer:

\begin{itemize}
  \item de baan is cirkelvormig;
  \item de baan ligt in het vlak van de evenaar;
  \item de satelliet beweegt in dezelfde draairichting als de aarde;
  \item de omlooptijd komt overeen met de rotatieperiode van de aarde
        ten opzichte van de sterren.
\end{itemize}

Uit

\[
T^2=\frac{4\pi^2}{GM_A}r^3
\]

volgt:

\[
\boxed{
r
=
\sqrt[3]{\frac{GM_AT^2}{4\pi^2}}.
}
\]

Dit is de afstand tot het aardmiddelpunt.

De hoogte boven het aardoppervlak is:

\[
\boxed{
h=r-R_A.
}
\]


\begin{oefening}[Geostationaire baan]{\oefU\ESS}

Neem voor een eerste berekening:

\[
T=23\,\mathrm h\,56\,\mathrm{min},
\]

\[
M_A=5{,}97\cdot10^{24}\,\mathrm{kg},
\qquad
R_A=6{,}37\cdot10^6\,\mathrm m.
\]

\begin{question}
Zet de periode om naar SI-eenheden.
\end{question}

\begin{question}
Bereken de vereiste baanstraal.
\end{question}

\begin{question}
Bereken de hoogte boven het aardoppervlak.
\end{question}

\begin{question}
Bereken de baansnelheid.
\end{question}

\begin{question}
Waarom zou een satelliet met dezelfde periode maar in een polaire baan
niet geostationair zijn?
\end{question}

\end{oefening}


% ============================================================
% 11 GEWICHTLOOSHEID
% ============================================================

\FVDsection{Gewichtloosheid is geen afwezigheid van gravitatie}

Astronauten in een ruimtestation lijken gewichtloos. Toch is het
gravitatieveld in een lage aardbaan nog aanzienlijk.

De verklaring is dat ruimtestation en astronaut samen vrijwel dezelfde
vrije val uitvoeren.

Als een astronaut los in het station zweeft, oefent de vloer geen
ondersteunende normaalkracht op de astronaut uit.

Dat is wat het gevoel van gewichtloosheid veroorzaakt.

\begin{opmerking}
De term ``gewicht'' wordt in verschillende contexten verschillend
gebruikt. In dit hoofdstuk is het belangrijk onderscheid te maken tussen
de gravitatiekracht $F_g$ en de contactkracht die bijvoorbeeld een
weegschaal meet.
\end{opmerking}


\begin{oefening}[Wat meet een weegschaal?]{\oefU\ESS}

Vergelijk een persoon:

\begin{enumerate}
  \item stilstaand op de aarde;
  \item in een lift in vrije val;
  \item zwevend in een ruimtestation in een cirkelbaan.
\end{enumerate}

\begin{question}
Is de gravitatiekracht in elk van de drie situaties noodzakelijk nul?
\end{question}

\begin{question}
Wanneer kan de normaalkracht van een weegschaal nul zijn?
\end{question}

\begin{question}
Leg uit waarom ``gewichtloos'' niet hetzelfde hoeft te betekenen als
``zonder gravitatiekracht''.
\end{question}

\end{oefening}


% ============================================================
% 12 MODELKRITIEK
% ============================================================

\FVDsection{Hoe goed is ons satellietmodel?}

We hebben tot nu toe vaak aangenomen dat:

\begin{itemize}
  \item de centrale massa bolsymmetrisch is;
  \item de satelliet als puntmassa kan worden behandeld;
  \item de baan perfect cirkelvormig is;
  \item alleen de gravitatiekracht van één centrale massa belangrijk is;
  \item luchtweerstand verwaarloosbaar is.
\end{itemize}

Echte banen kunnen elliptisch zijn en worden beïnvloed door andere
hemellichamen, de niet-perfecte bolsymmetrie van de aarde en, in lage
banen, een kleine atmosferische weerstand.

Een model is dus geen kopie van de werkelijkheid. Het is een
doelgerichte vereenvoudiging.


\begin{oefening}[Model of werkelijkheid?]{\oefG}

Een zeer lage satellietbaan wordt met het ideale cirkelbaanmodel
beschreven.

\begin{question}
Welke kracht ontbreekt mogelijk in het model?
\end{question}

\begin{question}
Verwacht je dat de mechanische energie van de satelliet in werkelijkheid
perfect constant blijft wanneer die kracht niet verwaarloosbaar is?
Motiveer kwalitatief.
\end{question}

\begin{question}
Waarom kan een model toch nuttig zijn wanneer het niet alle details van
de werkelijkheid bevat?
\end{question}

\end{oefening}


% ============================================================
% DEEL B ELEKTRISCH VELD
% ============================================================

\FVDsection{Transfer: bewegen in een elektrisch veld}

Gravitatie is niet het enige veld dat een kracht kan veroorzaken.

Voor een lading $q$ in een elektrisch veld geldt:

\[
\boxed{
\vec F_E=q\vec E.
}
\]

Voor een positieve lading heeft $\vec F_E$ dezelfde richting en zin als
$\vec E$.

Voor een negatieve lading is de zin tegengesteld.


\FVDsubsection{Newton II toepassen}

Als de elektrische kracht de enige relevante kracht is:

\[
m\vec a=q\vec E.
\]

Dus:

\[
\boxed{
\vec a=\frac qm\vec E.
}
\]

In een \textbf{homogeen} elektrisch veld is $\vec E$ constant. Voor een
deeltje met constante massa en lading is dan ook $\vec a$ constant.

We herkennen dus opnieuw een EVRB, maar nu veroorzaakt door een
elektrische kracht.


\begin{oefening}[Positief of negatief?]{\oefB}

Een geladen deeltje wordt vanuit rust losgelaten in een homogeen
elektrisch veld dat naar rechts wijst.

\begin{question}
In welke richting versnelt een positief deeltje?
\end{question}

\begin{question}
In welke richting versnelt een negatief deeltje?
\end{question}

\begin{question}
Waarom moet het teken van $q$ behouden blijven in de vectoriële
vergelijking
\[
\vec a=\frac qm\vec E?
\]
\end{question}

\end{oefening}


% ============================================================
% 14 GRAVITATIE VS ELEKTRISCH
% ============================================================

\FVDsection{Gravitatieveld en elektrisch veld vergelijken}

De twee veldconcepten vertonen een opvallende structuur.

Voor gravitatie:

\[
\vec F_g=m\vec g.
\]

Voor elektriciteit:

\[
\vec F_E=q\vec E.
\]

In beide gevallen is er:

\[
\boxed{
\text{eigenschap van het deeltje}
\times
\text{veldsterkte}
=
\text{kracht}.
}
\]

Maar er zijn ook belangrijke verschillen.

Massa is in de klassieke mechanica positief en gravitatie tussen
gewone massa's is aantrekkend.

Elektrische lading kan positief of negatief zijn. Daardoor kan de
elektrische kracht in de zin van $\vec E$ of in de tegengestelde zin
werken.


\begin{oefening}[Analogie met grenzen]{\oefU\ESS}

Vergelijk:

\[
\vec F_g=m\vec g
\]

en

\[
\vec F_E=q\vec E.
\]

\begin{question}
Welke rol speelt $m$ in de eerste vergelijking?
\end{question}

\begin{question}
Welke rol speelt $q$ in de tweede?
\end{question}

\begin{question}
Waarom vallen verschillende massa's in hetzelfde ideale gravitatieveld
met dezelfde versnelling?
\end{question}

\begin{question}
Waarom hoeven twee deeltjes met verschillende $q/m$-verhouding in
hetzelfde elektrische veld niet dezelfde versnelling te hebben?
\end{question}

\begin{question}
Waarom is een analogie nuttig, maar nooit een bewijs dat twee
verschijnselen fysisch identiek zijn?
\end{question}

\end{oefening}


% ============================================================
% 15 ELEKTRISCHE WORP
% ============================================================

\FVDsection{Een geladen deeltje tussen evenwijdige platen}

Beschouw een positief geladen deeltje dat horizontaal een gebied met
een homogeen verticaal elektrisch veld binnenkomt.

Kies:

\[
\vec v_0=v_0\vec e_x
\]

en een elektrisch veld in de positieve $y$-richting.

Dan:

\[
F_x=0
\]

en dus:

\[
a_x=0.
\]

Horizontaal:

\[
\boxed{
x(t)=x_0+v_0t.
}
\]

Verticaal:

\[
F_y=qE
\]

en:

\[
a_y=\frac{qE}{m}.
\]

Als $v_{y,0}=0$:

\[
\boxed{
y(t)
=
y_0+\frac12\frac{qE}{m}t^2.
}
\]

De beweging is dus de samenstelling van:

\begin{itemize}
  \item een eenparige horizontale beweging;
  \item een eenparig versnelde verticale beweging.
\end{itemize}

Dat is wiskundig dezelfde structuur als de horizontale worp in een
homogeen gravitatieveld.


\FVDsubsection{De baanvergelijking}

Uit:

\[
x-x_0=v_0t
\]

volgt:

\[
t=\frac{x-x_0}{v_0}.
\]

Invullen in de verticale bewegingsvergelijking geeft:

\[
y-y_0
=
\frac12\frac{qE}{m}
\left(\frac{x-x_0}{v_0}\right)^2.
\]

Dus:

\[
\boxed{
y(x)
=
y_0+
\frac{qE}{2mv_0^2}(x-x_0)^2.
}
\]

Binnen het gebied met homogeen elektrisch veld is de baan dus
parabolisch.


\begin{oefening}[Elektrische afbuiging]{\oefU}

Een positief geladen deeltje komt horizontaal met snelheid $v_0$ een
homogeen elektrisch veld binnen. Het veld wijst verticaal omhoog.

\begin{question}
Maak een krachten- en bewegingsanalyse per component.
\end{question}

\begin{question}
Leid $a_y=qE/m$ af.
\end{question}

\begin{question}
Stel $x(t)$ en $y(t)$ op voor $x_0=y_0=0$ en $v_{y,0}=0$.
\end{question}

\begin{question}
Elimineer $t$ en leid de baanvergelijking af.
\end{question}

\begin{question}
Wat verandert er aan de baan wanneer $q$ negatief is?
\end{question}

\begin{question}
Wat gebeurt er met de afbuiging wanneer $v_0$ groter wordt?
Interpreteer de factor $1/v_0^2$.
\end{question}

\end{oefening}


% ============================================================
% DEEL C MAGNETISCH VELD
% ============================================================

\FVDsection{Bewegen in een magnetisch veld}

Een bewegende elektrische lading kan in een magnetisch veld een
magnetische kracht ondervinden.

Voor de grootte geldt:

\[
\boxed{
F_B=|q|vB\sin\alpha,
}
\]

waar $\alpha$ de hoek is tussen $\vec v$ en $\vec B$.

De magnetische kracht staat loodrecht op zowel de snelheid als het
magnetisch veld.


\FVDsubsection{Drie bijzondere gevallen}

Als:

\[
\alpha=0^\circ
\quad\text{of}\quad
180^\circ,
\]

dan:

\[
F_B=0.
\]

Als:

\[
\alpha=90^\circ,
\]

dan is de kracht maximaal:

\[
\boxed{
F_B=|q|vB.
}
\]

Het geval $\vec v\perp\vec B$ is bijzonder interessant, omdat de kracht
dan voortdurend loodrecht op de snelheid kan staan.


\begin{oefening}[Eerst richting, dan grootte]{\oefB}

Een geladen deeltje beweegt in een homogeen magnetisch veld.

\begin{question}
Hoe groot is de magnetische kracht als het deeltje evenwijdig aan
$\vec B$ beweegt?
\end{question}

\begin{question}
Hoe groot is de magnetische kracht als het deeltje loodrecht op
$\vec B$ beweegt?
\end{question}

\begin{question}
Waarom kan een stilstaande lading door een zuiver magnetisch veld niet
worden versneld?
\end{question}

\end{oefening}


% ============================================================
% 17 MAGNETISCHE CIRKELBAAN
% ============================================================

\FVDsection{Wanneer magnetisme een cirkelbaan maakt}

Beschouw een geladen deeltje dat met snelheid $\vec v$ loodrecht op een
homogeen magnetisch veld $\vec B$ beweegt.

Dan is:

\[
F_B=|q|vB.
\]

De kracht staat op elk moment loodrecht op $\vec v$. Ze verandert dus
de richting van de snelheidsvector.

Wanneer het magnetische veld homogeen is en er geen andere relevante
krachten zijn, levert de magnetische kracht de radiale resultante:

\[
F_B=ma_c.
\]

Dus:

\[
|q|vB
=
m\frac{v^2}{r}.
\]

Voor $v\neq0$:

\[
|q|B
=
m\frac vr.
\]

Daaruit:

\[
\boxed{
r=\frac{mv}{|q|B}.
}
\]

Dit is een krachtige relatie: de kromming van de baan bevat informatie
over massa, lading, snelheid en magnetische veldsterkte.


\FVDsubsection{Wat vertelt de formule?}

Als de andere grootheden constant blijven:

\[
r\sim m,
\]

\[
r\sim v,
\]

\[
r\sim\frac1{|q|},
\]

en:

\[
r\sim\frac1B.
\]

Een groter $B$ buigt de baan dus sterker af en geeft een kleinere
straal.

Een sneller deeltje krijgt in hetzelfde veld een grotere baanstraal.


\begin{oefening}[Verbanden analyseren]{\oefU}

Een geladen deeltje beweegt loodrecht op een homogeen magnetisch veld.

\begin{question}
Wat gebeurt er met $r$ als $v$ verdubbelt?
\end{question}

\begin{question}
Wat gebeurt er met $r$ als $B$ verdubbelt?
\end{question}

\begin{question}
Wat gebeurt er met $r$ als zowel $m$ als $|q|$ verdubbelen?
\end{question}

\begin{question}
Waarom moet je bij elk antwoord vermelden welke andere grootheden
constant blijven?
\end{question}

\end{oefening}


% ============================================================
% 18 SNELHEIDSGROOTTE
% ============================================================

\FVDsection{Kan een magnetische kracht de snelheidsgrootte veranderen?}

Bij $\vec v\perp\vec B$ staat de magnetische kracht loodrecht op de
snelheid.

De kracht verandert daardoor de richting van $\vec v$, maar niet de
snelheidsgrootte.

Dit is precies wat nodig is voor een eenparige cirkelbeweging:

\[
|\vec v|=\text{constant},
\qquad
\vec v\neq\text{constant}.
\]

De versnelling is niet nul, want de richting van $\vec v$ verandert
voortdurend.

\begin{oefening}[Constante snelheid en toch versnellen]{\oefU}

Een leerling zegt:

\begin{quote}
``Als het magnetisch veld de snelheid niet groter maakt, kan het deeltje
ook geen versnelling hebben.''
\end{quote}

\begin{question}
Welke verwarring maakt de leerling?
\end{question}

\begin{question}
Welke grootheid blijft constant in een ideale magnetische cirkelbaan?
\end{question}

\begin{question}
Welke grootheid verandert wel voortdurend?
\end{question}

\begin{question}
Leg de link met de eenparige cirkelbeweging uit H5.
\end{question}

\end{oefening}


% ============================================================
% 19 PERIODE MAGNETISCHE BAAN
% ============================================================

\FVDsection{Verdieping: de omlooptijd in een magnetisch veld}

Voor een cirkelbeweging:

\[
T=\frac{2\pi r}{v}.
\]

Voor de magnetische cirkelbaan:

\[
r=\frac{mv}{|q|B}.
\]

Invullen geeft:

\[
T
=
\frac{2\pi}{v}
\frac{mv}{|q|B}.
\]

De snelheid valt weg:

\[
\boxed{
T=\frac{2\pi m}{|q|B}.
}
\]

En dus:

\[
\boxed{
\omega=\frac{|q|B}{m}.
}
\]

Binnen dit niet-relativistische ideale model hangt de omlooptijd dus
niet af van de snelheid.


\begin{oefening}[Een verrassend resultaat]{\oefG}

Twee identieke positief geladen deeltjes bewegen loodrecht op hetzelfde
homogene magnetische veld.

Voor de snelheden geldt:

\[
v_2=3v_1.
\]

\begin{question}
Bepaal $r_2/r_1$.
\end{question}

\begin{question}
Bepaal $T_2/T_1$.
\end{question}

\begin{question}
Hoe kan het tweede deeltje een grotere cirkel afleggen in dezelfde
omlooptijd?
\end{question}

\begin{question}
Welke aannames liggen aan deze conclusie ten grondslag?
\end{question}

\end{oefening}


% ============================================================
% 20 TEKEN VAN LADING
% ============================================================

\FVDsection{De zin van de afbuiging}

De formule

\[
r=\frac{mv}{|q|B}
\]

bevat $|q|$, omdat de straal positief is.

Maar het \textbf{teken} van de lading is wel essentieel voor de zin van
de magnetische kracht.

Een positieve en een negatieve lading met dezelfde massa,
snelheidsgrootte en $|q|$ krijgen in hetzelfde magnetische veld een
baan met dezelfde straal, maar buigen in tegengestelde zin af.

\begin{oefening}[Zelfde straal, andere zin]{\oefU}

Een positief en een negatief deeltje hebben dezelfde massa, dezelfde
$|q|$ en dezelfde beginsnelheid. Ze komen loodrecht hetzelfde homogene
magnetische veld binnen.

\begin{question}
Vergelijk de grootte van de magnetische krachten.
\end{question}

\begin{question}
Vergelijk de baanstralen.
\end{question}

\begin{question}
Vergelijk de zin van de afbuiging.
\end{question}

\begin{question}
Waarom mag je in de straalformule $|q|$ gebruiken, maar niet vergeten
dat het teken van $q$ fysische informatie bevat?
\end{question}

\end{oefening}


% ============================================================
% 21 q/m
% ============================================================

\FVDsection{De baan als meetinstrument}

Uit:

\[
r=\frac{mv}{|q|B}
\]

volgt:

\[
\boxed{
\frac{m}{|q|}
=
\frac{Br}{v}.
}
\]

Een gemeten baanstraal kan dus, als $B$ en $v$ bekend zijn, informatie
geven over de verhouding tussen massa en ladingsgrootte.

Dat principe ligt aan de basis van technieken waarbij geladen deeltjes
worden gescheiden of geïdentificeerd op basis van hun beweging in
velden.

\begin{oefening}[Onbekend deeltje]{\oefG}

Een positief geladen deeltje beweegt loodrecht op een homogeen
magnetisch veld. De grootheden $B$, $v$ en $r$ worden experimenteel
bepaald.

\begin{question}
Leid een formule af voor $m/|q|$.
\end{question}

\begin{question}
Welke metingen zijn nodig?
\end{question}

\begin{question}
Als twee deeltjes dezelfde snelheid hebben en in hetzelfde veld een
verschillende straal beschrijven, wat kan je dan over hun
$m/|q|$-verhouding besluiten?
\end{question}

\begin{question}
Waarom kan de baanstraal alleen niet altijd volstaan om een deeltje
uniek te identificeren?
\end{question}

\end{oefening}


% ============================================================
% 22 GEKRUISTE VELDEN
% ============================================================

\FVDsection{Gevorderd: elektrische en magnetische kracht samen}

Velden kunnen tegelijk aanwezig zijn.

Stel dat een geladen deeltje door een gebied beweegt waarin een
elektrisch en een magnetisch veld zo georiënteerd zijn dat de
elektrische en magnetische kracht tegengestelde zin hebben.

Voor een rechtlijnige beweging zonder afbuiging moet:

\[
\sum\vec F=\vec0.
\]

In grootte:

\[
F_E=F_B.
\]

Voor loodrechte snelheid en magnetisch veld:

\[
|q|E=|q|vB.
\]

Dus:

\[
\boxed{
v=\frac EB.
}
\]

Alleen deeltjes met deze snelheid gaan in het ideale model rechtdoor.


\begin{oefening}[Snelheidsselector]{\oefG}

Een bundel geladen deeltjes beweegt door gekruiste homogene elektrische
en magnetische velden.

\begin{question}
Leid de voorwaarde
\[
v=\frac EB
\]
af voor een niet-afgebogen deeltje.
\end{question}

\begin{question}
Waarom verdwijnt $q$ uit deze formule?
\end{question}

\begin{question}
Wat gebeurt er kwalitatief met een deeltje dat sneller beweegt dan
$E/B$, als de twee krachten bij $v=E/B$ tegengesteld zijn?
\end{question}

\begin{question}
Waarom is dit een voorbeeld van technologie die een fysisch model als
selectiemechanisme gebruikt?
\end{question}

\end{oefening}


% ============================================================
% 23 GROTE SYNTHESE
% ============================================================

\FVDsection{Drie velden, drie bewegingspatronen}

We kunnen nu verschillende situaties naast elkaar plaatsen.

\[
\boxed{
\vec F_g=m\vec g
}
\]

in een ongeveer homogeen gravitatieveld geeft:

\[
\vec a=\vec g.
\]

Een constante veldvector geeft dus een constante versnelling.

Voor een homogeen elektrisch veld:

\[
\boxed{
\vec F_E=q\vec E
}
\]

en:

\[
\boxed{
\vec a=\frac qm\vec E.
}
\]

Ook hier kan de versnelling constant zijn.

Voor een magnetisch veld:

\[
\boxed{
F_B=|q|vB\sin\alpha.
}
\]

Bij $\vec v\perp\vec B$ staat de kracht voortdurend loodrecht op de
snelheid en kan een cirkelbeweging ontstaan.

Bij een radiaal gravitatieveld rond een bolmassa:

\[
g(r)=\frac{GM}{r^2},
\]

wijst de kracht naar het centrum. Bij de juiste tangentiële snelheid
kan ook daar een cirkelbaan ontstaan.


\FVDsubsection{De structuur achter de formules}

Het belangrijkste inzicht van dit hoofdstuk is niet één afzonderlijke
formule, maar een strategie:

\[
\boxed{
\begin{array}{c}
\text{1. Welk veld is aanwezig?}\\[2mm]
\downarrow\\[2mm]
\text{2. Welke kracht werkt op het gekozen systeem?}\\[2mm]
\downarrow\\[2mm]
\text{3. Wat zegt }\sum\vec F=m\vec a\text{?}\\[2mm]
\downarrow\\[2mm]
\text{4. Welk bewegingsmodel past bij }\vec a\text{?}\\[2mm]
\downarrow\\[2mm]
\text{5. Wat voorspelt het model?}
\end{array}
}
\]


\begin{oefening}[Herken het bewegingsmodel]{\oefU\ESS}

Koppel elk geval aan het meest geschikte ideale bewegingsmodel en
motiveer vanuit de kracht.

\begin{enumerate}
  \item Een geladen deeltje start uit rust in een homogeen elektrisch veld.
  \item Een geladen deeltje beweegt evenwijdig aan een homogeen magnetisch
        veld en er werken geen andere krachten.
  \item Een geladen deeltje beweegt loodrecht op een homogeen magnetisch veld.
  \item Een satelliet beweegt met geschikte tangentiële snelheid rond een
        bolsymmetrische planeet.
  \item Een voorwerp wordt horizontaal geworpen nabij het aardoppervlak.
\end{enumerate}

Gebruik waar passend de termen EVRB, ERB, horizontale worp en ECB.

\end{oefening}


% ============================================================
% 24 MATHEMATISEREN
% ============================================================

\FVDsection{Een veldprobleem systematisch oplossen}

Gebruik bij nieuwe problemen niet meteen een formule. Werk vanuit de
fysica.

\begin{enumerate}
  \item \textbf{Kies het systeem.}
        Is het een satelliet, ion, elektron, planeet of ander object?
  \item \textbf{Identificeer het veld.}
        Gravitatie, elektrisch, magnetisch of een combinatie?
  \item \textbf{Bepaal de kracht.}
        Bijvoorbeeld
        \[
        F_g=G\frac{Mm}{r^2},
        \quad
        \vec F_E=q\vec E,
        \quad
        F_B=|q|vB\sin\alpha.
        \]
  \item \textbf{Gebruik Newton II.}
        \[
        \sum\vec F=m\vec a.
        \]
  \item \textbf{Herken het bewegingsmodel.}
        Is $\vec a$ constant? Staat $\vec a$ loodrecht op $\vec v$?
        Is de versnelling radiaal?
  \item \textbf{Voeg de juiste kinematica toe.}
        Bijvoorbeeld
        \[
        a_c=\frac{v^2}{r}
        \]
        of de bewegingsvergelijkingen bij constante versnelling.
  \item \textbf{Los symbolisch op waar dat zinvol is.}
  \item \textbf{Controleer dimensies, richting en grootteorde.}
  \item \textbf{Interpreteer het resultaat fysisch.}
\end{enumerate}


% ============================================================
% 25 GEMENGDE OEFENINGEN
% ============================================================

\FVDsection{Geïntegreerde oefeningen}

\begin{oefening}[Van $g$ naar een satellietbaan]{\oefU\ESS}

Op afstand $r$ van het middelpunt van een planeet wordt de
gravitatieveldsterkte gemeten als $g_r$.

Een satelliet beweegt op diezelfde afstand in een ideale cirkelbaan.

\begin{question}
Toon zonder de massa van de planeet expliciet te gebruiken dat:
\[
\boxed{
v=\sqrt{g_rr}.
}
\]
\end{question}

\begin{question}
Leid vervolgens af:
\[
\boxed{
T=2\pi\sqrt{\frac r{g_r}}.
}
\]
\end{question}

\begin{question}
Waarom is dit een mooi voorbeeld van het combineren van twee eerder
geleerde modellen?
\end{question}

\end{oefening}


\begin{oefening}[Een onbekende planeet]{\oefG}

Een maan beweegt in een vrijwel cirkelvormige baan rond een onbekende
planeet. De baanstraal en omlooptijd zijn gemeten als $r$ en $T$.

\begin{question}
Leid af:
\[
\boxed{
M=\frac{4\pi^2r^3}{GT^2}.
}
\]
\end{question}

\begin{question}
Welke massa verschijnt niet in het eindresultaat?
\end{question}

\begin{question}
Leg uit waarom de baan van een maan informatie bevat over de massa van
de planeet.
\end{question}

\begin{question}
Welke experimentele grootheden bepalen vooral de onzekerheid op de
berekende massa?
\end{question}

\end{oefening}


\begin{oefening}[Elektrisch of gravitatieel?]{\oefU}

Een deeltje beweegt door een gebied waarin een constante kracht naar
beneden werkt. De baan is parabolisch.

\begin{question}
Kan je alleen uit de vorm van de baan besluiten dat de kracht
gravitatieel is? Motiveer.
\end{question}

\begin{question}
Geef een elektrische situatie die wiskundig dezelfde soort baan kan
opleveren.
\end{question}

\begin{question}
Welke bijkomende informatie zou helpen om beide modellen te
onderscheiden?
\end{question}

\end{oefening}


\begin{oefening}[Twee ionen in hetzelfde magneetveld]{\oefG}

Twee positief geladen ionen komen met dezelfde snelheid loodrecht
hetzelfde homogene magnetische veld binnen.

Voor de eerste geldt:

\[
m_1=m,\qquad q_1=e.
\]

Voor de tweede:

\[
m_2=4m,\qquad q_2=2e.
\]

\begin{question}
Bepaal $r_2/r_1$.
\end{question}

\begin{question}
Bepaal $T_2/T_1$.
\end{question}

\begin{question}
Welke van beide banen is sterker gekromd?
\end{question}

\begin{question}
Welke verhouding $m/|q|$ wordt door de baan zichtbaar?
\end{question}

\end{oefening}


\begin{oefening}[Welke uitspraak klopt?]{\oefU\ESS}

Beoordeel telkens de uitspraak en verbeter ze indien nodig.

\begin{enumerate}
  \item ``Een satelliet is gewichtloos omdat $g=0$.''
  \item ``Centripetale kracht is een extra kracht die naast gravitatie
        op een satelliet werkt.''
  \item ``Een hogere cirkelbaan vraagt een grotere baansnelheid.''
  \item ``Een homogeen elektrisch veld kan een constante versnelling
        veroorzaken.''
  \item ``Een magnetische kracht kan bij $\vec v\perp\vec B$ de richting
        van de snelheid veranderen zonder de snelheidsgrootte te veranderen.''
  \item ``Als twee deeltjes in hetzelfde magneetveld dezelfde baanstraal
        hebben, moeten ze dezelfde massa hebben.''
\end{enumerate}

Motiveer elk antwoord met een fysisch verband en niet alleen met woorden.

\end{oefening}


% ============================================================
% 26 ICT EN DATA
% ============================================================

\FVDsection{Onderzoeken met data en ICT}

Dit hoofdstuk leent zich sterk tot modelleren en gegevensanalyse.

Voor satellieten voorspelt het model:

\[
T^2=\frac{4\pi^2}{GM}r^3.
\]

Met echte of gegeven satellietgegevens kan je:

\begin{enumerate}
  \item $r$ en $T$ verzamelen;
  \item $r^3$ en $T^2$ berekenen;
  \item een spreidingsdiagram van $T^2$ tegen $r^3$ maken;
  \item een lineaire regressie uitvoeren;
  \item de helling interpreteren;
  \item daaruit $M$ schatten;
  \item afwijkingen en meetonzekerheid bespreken.
\end{enumerate}

Zo wordt een theoretische afleiding een toetsbaar model.


\begin{oefening}[Onderzoeksopdracht: Kepler uit data]{\oefU\ESS}

Gebruik een dataset met minstens vier manen of satellieten rond
hetzelfde centrale hemellichaam.

\begin{question}
Maak een tabel met $r$, $T$, $r^3$ en $T^2$.
\end{question}

\begin{question}
Maak met ICT een grafiek van $T^2$ tegen $r^3$.
\end{question}

\begin{question}
Voer een geschikte regressie uit.
\end{question}

\begin{question}
Interpreteer de helling fysisch.
\end{question}

\begin{question}
Gebruik de helling om de massa van het centrale hemellichaam te
schatten.
\end{question}

\begin{question}
Vergelijk met een referentiewaarde en bespreek mogelijke oorzaken van
een verschil.
\end{question}

\end{oefening}


% ============================================================
% 27 VERDIEPING ALGEMENE MAGNETISCHE BEWEGING
% ============================================================

\FVDsection{Gevorderd: wat als $\vec v$ niet loodrecht op $\vec B$ staat?}

Tot nu toe namen we bij de magnetische cirkelbaan:

\[
\vec v\perp\vec B.
\]

Bij een algemene invalshoek kunnen we de snelheid ontbinden in:

\[
\vec v=\vec v_\parallel+\vec v_\perp.
\]

De component $v_\parallel$ veroorzaakt geen magnetische kracht:

\[
F_{B,\parallel}=0.
\]

De loodrechte component $v_\perp$ veroorzaakt wel een magnetische
afbuiging.

De combinatie van:

\begin{itemize}
  \item een eenparige beweging langs de veldrichting;
  \item een cirkelbeweging loodrecht op de veldrichting
\end{itemize}

levert een schroefvormige of helicale baan.

De straal van de cirkelcomponent is:

\[
\boxed{
r=\frac{mv_\perp}{|q|B}.
}
\]

Dit is een mooi voorbeeld van het onafhankelijk analyseren van
vectorcomponenten.


\begin{oefening}[Helix in een magnetisch veld]{\oefG}

Een geladen deeltje komt onder een hoek in een homogeen magnetisch veld.

\begin{question}
Ontbind $\vec v$ in $v_\parallel$ en $v_\perp$.
\end{question}

\begin{question}
Welke component bepaalt de straal van de baan?
\end{question}

\begin{question}
Welke component bepaalt hoe snel het deeltje langs de veldrichting
opschuift?
\end{question}

\begin{question}
Wat voor baan ontstaat als $v_\parallel=0$?
\end{question}

\begin{question}
Wat voor baan ontstaat als $v_\perp=0$?
\end{question}

\end{oefening}


% ============================================================
% 28 WISKUNDIGE VERDIEPING
% ============================================================

\FVDsection{Wiskundige verdieping: machtswetten herkennen}

Verschillende relaties uit dit hoofdstuk zijn machtsverbanden.

Voor gravitatie:

\[
g(r)=GM\,r^{-2}.
\]

Voor een satelliet:

\[
v(r)=\sqrt{GM}\,r^{-1/2},
\]

en:

\[
T(r)=\frac{2\pi}{\sqrt{GM}}r^{3/2}.
\]

De exponent vertelt onmiddellijk hoe gevoelig een grootheid is voor een
verandering van $r$.

Bijvoorbeeld:

\[
r\rightarrow kr
\]

geeft:

\[
g\rightarrow k^{-2}g,
\]

\[
v\rightarrow k^{-1/2}v,
\]

\[
T\rightarrow k^{3/2}T.
\]


\begin{oefening}[Eén verandering, drie gevolgen]{\oefG}

Een satelliet verhuist in een ideaal model van een cirkelbaan met
straal $r$ naar een cirkelbaan met straal $9r$ rond dezelfde planeet.

Vergelijk voor beide banen:

\begin{question}
de gravitatieveldsterkte;
\end{question}

\begin{question}
de baansnelheid;
\end{question}

\begin{question}
de omlooptijd.
\end{question}

Los op met de machtsverbanden, zonder de formules telkens volledig
numeriek in te vullen.

\end{oefening}


% ============================================================
% 29 WAT NIET IN DIT HOOFDSTUK
% ============================================================

\FVDsection{Vooruitblik: velden en energie}

Een gravitatieveld kan ook energetisch beschreven worden. In je eerdere
leerstof ontmoette je mogelijk al gravitatiepotentiaal en
gravitatiepotentiële energie.

Voor grote afstanden is de eenvoudige nabij-aardeformule

\[
E_{p}=mgh
\]

niet algemeen genoeg. Dan verschijnt een afstandsafhankelijke
gravitatie-energie.

We bouwen die energiebeschrijving hier bewust nog niet volledig uit.
Ze past inhoudelijk beter nadat arbeid en energie in deze cursus
systematisch behandeld zijn.

Dan kunnen we onder andere onderzoeken:

\begin{itemize}
  \item gravitatiepotentiële energie op grote afstand;
  \item arbeid in een gravitatieveld;
  \item totale mechanische energie van banen;
  \item ontsnappingssnelheid.
\end{itemize}

Zo vermijden we dat kracht-, bewegings- en energiemodellen door elkaar
worden geïntroduceerd voordat hun eigen betekenis helder is.


% ============================================================
% 30 SAMENVATTING
% ============================================================

\FVDsection{Samenvatting}

Voor gravitatie:

\[
\boxed{
F_g=G\frac{Mm}{r^2}
}
\]

en:

\[
\boxed{
g(r)=G\frac{M}{r^2}.
}
\]

Voor een ideale cirkelvormige satellietbaan:

\[
G\frac{Mm}{r^2}
=
m\frac{v^2}{r}.
\]

Daaruit:

\[
\boxed{
v_{\mathrm{baan}}
=
\sqrt{\frac{GM}{r}}
}
\]

en:

\[
\boxed{
T
=
2\pi\sqrt{\frac{r^3}{GM}}.
}
\]

Dus voor dezelfde centrale massa:

\[
\boxed{
T^2\sim r^3.
}
\]

Voor een elektrisch veld:

\[
\boxed{
\vec F_E=q\vec E
}
\]

en in een homogeen elektrisch veld:

\[
\boxed{
\vec a=\frac qm\vec E.
}
\]

Voor een bewegende lading in een magnetisch veld:

\[
\boxed{
F_B=|q|vB\sin\alpha.
}
\]

Bij $\vec v\perp\vec B$ en zonder andere relevante krachten:

\[
|q|vB=m\frac{v^2}{r}
\]

en dus:

\[
\boxed{
r=\frac{mv}{|q|B}.
}
\]

Als verdieping:

\[
\boxed{
T=\frac{2\pi m}{|q|B}.
}
\]

Doorheen al deze situaties blijft dezelfde analysekern gelden:

\[
\boxed{
\text{veld}
\rightarrow
\text{kracht}
\rightarrow
\text{Newton II}
\rightarrow
\text{versnelling}
\rightarrow
\text{beweging}.
}
\]



% ============================================================
% 32 SLOT
% ============================================================

\FVDsection{Van krachten naar energie}

We begonnen dit thema met de vraag hoe we beweging kunnen beschrijven.
Daarna onderzochten we waardoor beweging verandert.

In dit hoofdstuk zagen we hoe ver die ideeën reiken. Dezelfde wetten
verklaren zowel bewegingen op astronomische schaal als de afbuiging van
geladen deeltjes.

Maar krachten zijn niet de enige taal waarmee we beweging kunnen
analyseren.

Soms is het veel efficiënter om niet elk moment van de beweging te
volgen, maar te onderzoeken hoeveel energie wordt overgedragen of
omgezet.

Daarmee ontstaat de volgende stap:

\[
\boxed{
\text{kracht en beweging}
\longrightarrow
\text{arbeid en energie}.
}
\]

\end{document}
